\section{总结与延伸}

原文作者 Mr Panda 在大半年的 Claude Code 实战中提炼出一套清晰的使用哲学：

\begin{importantbox}{Claude Code 核心使用哲学}
它是一个极其能干的协作者，但不是自动驾驶。方向盘始终在你手里。工具的价值，永远取决于使用它的人。
\end{importantbox}

按重要性排序的六条建议：

\begin{enumerate}[leftmargin=1.5em]
  \item \textbf{写好 CLAUDE.md}——一次投入，每次对话都受益
  \item \textbf{给精确的指令}——目标 + 约束 + 验证标准
  \item \textbf{用 Plan Mode}——复杂任务先规划再动手
  \item \textbf{管理上下文}——\texttt{/clear}、\texttt{/compact}、子 Agent
  \item \textbf{控制权限}——deny 危险操作，allow 常用命令
  \item \textbf{频繁 commit}——保留回退能力
\end{enumerate}

\subsection{拓展阅读}

\begin{itemize}[leftmargin=1.5em]
  \item Claude Code 官方文档：\url{https://docs.anthropic.com/en/docs/claude-code}
  \item Model Context Protocol 规范：\url{https://modelcontextprotocol.io}
  \item 原文链接：\url{https://x.com/PandaTalk8/status/2035975664730575325}
\end{itemize}

\end{document}
